\documentclass[11pt]{article}
\usepackage{amsmath,amssymb}
\usepackage{amsthm}
\usepackage{mathtools}
\usepackage[margin=1in]{geometry}
\usepackage{booktabs}
\usepackage{array}
\usepackage{hyperref}
\usepackage{xcolor}
\hypersetup{colorlinks=true, linkcolor=blue!50!black, urlcolor=blue!50!black,
            citecolor=blue!50!black}

\newcommand{\code}[1]{\texttt{#1}}
\newcommand{\Gm}{\mathbb{G}_m}
\newcommand{\Ga}{\mathbb{G}_a}
\newcommand{\CC}{\mathbb{C}}
\newcommand{\ZZ}{\mathbb{Z}}
\newcommand{\QQ}{\mathbb{Q}}
\newcommand{\gl}{\mathfrak{gl}}
\newcommand{\dprem}{\operatorname{d-prem}}

\theoremstyle{plain}
\newtheorem{proposition}{Proposition}
\newtheorem{lemma}[proposition]{Lemma}
\theoremstyle{definition}
\newtheorem{criterion}[proposition]{Criterion}
\theoremstyle{remark}
\newtheorem{remark}[proposition]{Remark}

\title{Detecting and spending symmetry in a parametric ansatz\\
       {\large three tiers of algorithm, and a per-cell normalization criterion}}
\author{}
\date{6 September 2026}

\begin{document}
\maketitle

\begin{abstract}
\noindent
A companion note to \emph{A New Method of Solving PDEs}, backing the
\S\,Scaling Symmetries subsection. The parameter count of an ansatz is the
dominant cost driver of the decomposition that follows it, so every dimension
of symmetry that can be pinned is worth a parameter and its constancy
equations. This note asks what can be \emph{detected algorithmically}, and
answers in three tiers: diagonal (scaling) symmetries are an integer kernel
computation and are complete for the torus; all \emph{linear} symmetries are a
linear solve in $\gl_N$ and additionally capture the unipotent directions that
a weight table is structurally blind to; general Lie point symmetries require
the classical determining system of Olver, solved by the same differential
elimination already in the pipeline. We then formulate a
\emph{per-cell} normalization criterion --- at the moment a decomposition
splits on $k \neq 0$, when may one set $k = 1$? --- give its exact condition,
show that the available symmetry is \emph{non-decreasing} as one descends the
decomposition tree (so per-cell normalization strictly dominates global
normalization), and identify precisely which symmetries the criterion misses.
The Tier-1 lattice is computed for the hydrogen / ansatz~5 system: the torus is
$2$-dimensional, the energy $E$ is provably unpinnable, and five of the
``obvious'' two-parameter normalizations --- including $a_0 = c_0 = 1$ ---
silently produce a $d$-fold cover of the answer.
\end{abstract}

\section{Why this matters: parameters are the cost}
\label{sec:cost}

The ansatz contributes a fixed number of differential equations, but each
parameter contributes a family of constancy equations --- one per (parameter,
coordinate) pair. For hydrogen / ansatz~5 that is $30$ constancy equations for
$10$ parameters, three apiece, against a fixed $9$ ansatz equations. Every
dimension of symmetry that can be pinned is therefore worth exactly one
parameter and its constancy equations, and the question of the present note is
which of those dimensions a machine can find without being told.

The stakes are not marginal. Table~\ref{tab:hydrogen5} gives the per-cell cost
of the membership run for hydrogen / ansatz~5 completed on 6~September 2026
(\code{hydrogen-5-membership.out}; \code{--pde hydrogen --ansatz 5 --locus
membership --basic}).

\begin{table}[htbp]
\centering
\small
\begin{tabular}{r l r r r r}
\toprule
cell & vanishing set & \#ineq & $\dprem$ (s) & GTZ (s) & primes \\
\midrule
 1 & (generic)                                 & 6 &  15.3 &    2.7 & 4 \\
 2 & $a_0$                                     & 3 &   9.8 &    0.4 & 4 \\
 3 & $v_4$                                     & 4 &   0.1 &    0.0 & 2 \\
 4 & $a_0, v_4$                                & 2 &   0.0 &    0.0 & 2 \\
 5 & $a_1$                                     & 3 &   0.9 &    0.1 & 1 \\
 6 & $a_1, v_4$                                & 2 &   0.0 &    0.0 & 1 \\
 7 & $a_0, a_1, v_4$                           & 4 &   1.8 &    0.1 & 2 \\
 8 & $a_0, a_1, b_0, v_4$                      & 2 &   0.4 &    0.0 & 2 \\
 9 & $a_0, a_1, b_1, v_4$                      & 2 &   0.0 &    0.0 & 1 \\
10 & $a_0, a_1, b_0, b_1, c_0, c_1, v_4$       & 2 &   0.0 &    0.0 & --- \\
11 & $\;\;\cdots, v_3, v_4$                    & 2 &   0.0 &    0.0 & --- \\
12 & $\;\;\cdots, v_2, v_3, v_4$               & 1 &   0.0 &    0.0 & --- \\
\textbf{13} & $\mathbf{a_0, a_1}$              & \textbf{7} & \textbf{476.4} & \textbf{1430.2} & 4 \\
14 & $a_0, a_1, b_0$                           & 4 & 167.9 &    1.3 & 3 \\
15 & $a_0, a_1, b_1$                           & 4 &   4.2 &    0.1 & 2 \\
16 & $a_0, a_1, b_0, b_1, c_0, c_1$            & 3 &   0.9 &    0.2 & 1 \\
\bottomrule
\end{tabular}
\caption{Per-cell cost, hydrogen / ansatz~5 membership run. The Thomas
decomposition itself took $495.2$~s and produced these $16$ cells; the
per-cell work totals $2112.8$~s. \textbf{Cell~13 alone is $90.2\,\%$ of the
per-cell work} and $73\,\%$ of the run.}
\label{tab:hydrogen5}
\end{table}

Two observations shape everything below. First, the cost is
\emph{overwhelmingly concentrated}: one cell out of sixteen carries nine tenths
of the work, and cells 13 and 14 together carry $98.2\,\%$. A normalization
that helps the cheap cells is worthless; a normalization that helps cell~13 is
the whole game. Second, cell~13 is not the generic cell. It is the stratum
$a_0 = a_1 = 0$, on which the coefficient $(a_0 + a_1 v)$ of $\Psi''$ vanishes
identically and the ansatz ODE \emph{degenerates from second order to first}.
It carries the most surviving inequations of any cell ($7$), which is to say the
most candidates for the criterion of \S\ref{sec:criterion}.

\section{Two kinds of symmetry, and why the distinction is not cosmetic}
\label{sec:twokinds}

Write $P$ for the parameter space of an ansatz. It carries an action of a group
$G$ that splits, as in \cite{ansatz25}, according to whether a transformation
moves the solution set or only its description:

\begin{itemize}
\item $\Gamma$, the \textbf{gauge group}: the kernel of
      $\text{parameters} \mapsto \text{solution}$. A $\Gamma$-orbit is a set of
      parameter points describing \emph{the same} family of solutions ---
      multiplying the ODE through by a constant is the archetype. Quotienting
      by $\Gamma$ is lossless and incurs no obligation.
\item $H$, the \textbf{form-preserving symmetries}: genuine point symmetries of
      the PDE that carry the ansatz family to itself --- rotation, scaling of
      the independent variable. These \emph{permute} solutions, so pinning one
      computes a slice whose answer must afterwards be spread back over the
      orbit.
\end{itemize}

Both are worth spending; they differ only in whether a transport step is owed
at the end. As will be seen in \S\ref{sec:criterion}, the weight vector itself
says which case one is in, so the distinction costs nothing to maintain.

\section{Tier 1: diagonal (scaling) symmetries}
\label{sec:tier1}

\subsection{The algorithm}

A scaling symmetry is a character of a torus, so the entire group is the kernel
of an integer matrix and no search is involved.

Assign an unknown weight $w_s \in \ZZ$ to every symbol $s$ --- independent
variables, differential indeterminates, parameters, and the equation-multiplier.
A derivative $\partial^j u / \partial x^j$ carries weight $w_u - j\,w_x$. A
polynomial $f = \sum_\alpha c_\alpha z^\alpha$ is \emph{$w$-quasi-homogeneous}
iff $\langle \alpha - \alpha_0, w\rangle = 0$ for every $\alpha \in
\operatorname{supp} f$ and one fixed base exponent $\alpha_0$. Stack those rows
over all generators to get a matrix $M$; the scaling symmetries are
$\ker M$ over $\QQ$, and a Hermite normal form yields a $\ZZ$-basis. The
computation is exact, complete for the torus, requires no Gr\"obner basis, and
is instantaneous at the sizes that arise here.

This is not a new construction. It is the algorithm of Hubert and Labahn
\cite{HubertLabahnISSAC,HubertLabahnFoCM}, who show that the Hermite normal
form of the weight matrix produces a change of variables separating the group
directions from a complete generating set of rational invariants. (It is
perhaps worth noting that this is the same Hubert of \cite{hubert,hubert2},
already cited in the paper for triangular-set decomposition.)

\subsection{The worked example of the \S\,Scaling Symmetries subsection}

For $a\Psi'' + b\Psi' + c\Psi = 0$, with $w_x$, $w_\Psi$ and the parameter
weights unknown, quasi-homogeneity of the single generator reads
\[
   w_a + w_\Psi - 2w_x \;=\; w_b + w_\Psi - w_x \;=\; w_c + w_\Psi ,
\]
i.e.\ $w_a = w_c + 2w_x$ and $w_b = w_c + w_x$. The weight $w_\Psi$ cancels
identically --- the equation is linear homogeneous in $\Psi$ --- so it acts
trivially on the parameters, and the effective torus on $(a,b,c)$ is
\textbf{$2$-dimensional}, spanned by $w_c$ and $w_x$:
\[
  a : (1,2), \qquad b : (1,1), \qquad c : (1,0)
  \qquad\text{in coordinates } (w_c, w_x).
\]
The two basis directions are exactly the two symmetries of the draft
subsection: $(w_c, w_x) = (1,0)$ gives $(w_a,w_b,w_c) = (1,1,1)$, the overall
multiplication by $\lambda$; and $(0,1)$ gives $(2,1,0)$, the scaling of the
independent variable with coefficients scaled by $\lambda^2, \lambda, 1$. The
rank being $2$, exactly two of the three constants may be pinned and one
genuine modulus survives --- the classical fact that a second-order linear ODE
has one invariant.

\subsection{Unimodularity, and why it is not pedantry}
\label{sec:unimodular}

Pinning $k$ parameters kills $(\Gm)^k$ cleanly \emph{iff the corresponding
$k \times k$ weight minor has determinant $\pm 1$}. If the determinant is $d$
with $|d| > 1$, the map $(\Gm)^k \to (\Gm)^k$ is a $d$-fold isogeny: the slice
meets each orbit $d$ times, and the decomposition returns $d$ copies of every
component. Non-vanishing of the minor is \emph{not} the condition;
unimodularity is.

The example above already exhibits the trap. Of the three two-parameter pins:
\[
  \det\begin{pmatrix}1&2\\1&1\end{pmatrix} = -1 \;\;(a{=}b{=}1)\ \checkmark
  \qquad
  \det\begin{pmatrix}1&1\\1&0\end{pmatrix} = -1 \;\;(b{=}c{=}1)\ \checkmark
\]
\[
  \det\begin{pmatrix}1&2\\1&0\end{pmatrix} = -2 \;\;(a{=}c{=}1)
  \qquad\longleftarrow\ \textbf{defective}
\]
So the most symmetric-looking normalization, $\Psi'' + b\Psi' + \Psi = 0$, is
a $2$-fold cover. The residual $\mu_2$ acts by $b \mapsto -b$, and indeed
$\Psi'' + b\Psi' + \Psi$ and $\Psi'' - b\Psi' + \Psi$ are carried to one
another by $x \mapsto -x$. The defect is real, it is invisible without the
determinant check, and it doubles the component count of everything
downstream.

\subsection{Result: the hydrogen / ansatz 5 lattice}
\label{sec:hydrogen}

Applying the algorithm to the full system of \S\ref{sec:cost} --- nine ansatz
equations, the Schr\"odinger PDE, and the two inequations --- yields a
$3$-dimensional weight lattice with coordinates $(\Theta, \mu, P)$, where
$\Theta$ is the ODE-multiplier direction, $\mu$ the scaling of the inner
variable $v$, and $P$ the scaling of $\Psi$. Since $P$ acts trivially on
parameters, the effective torus on parameter space is \textbf{$2$-dimensional}:

\begin{center}
\begin{tabular}{l c c c c c c c c}
\toprule
 & $v_1 \ldots v_4$ & $a_0$ & $a_1$ & $b_0$ & $b_1$ & $c_0$ & $c_1$ & $E$ \\
\midrule
$(\Theta, \mu)$ & $(0,1)$ & $(1,2)$ & $(1,1)$ & $(1,1)$ & $(1,0)$ & $(1,0)$ & $(1,-1)$ & $(0,0)$ \\
\bottomrule
\end{tabular}
\end{center}

Every generator was checked individually
(\code{symmetry-weights-check.py}, \S\ref{sec:checked}); the table reproduces
the general formula of \cite{ansatz25} --- an affine coefficient pair
$(x_0 + x_1 v)$ multiplying a jet monomial of derivative order $w_j$ carries
$\mu$-weights $w_j$ and $w_j - 1$ --- which is an independent confirmation.

Three consequences, all of them useful:

\begin{enumerate}
\item \textbf{$E$ has weight $(0,0)$ and is therefore never pinnable.} This is
      forced, not chosen: the PDE term $\Psi$ (from the Coulomb $2\Psi/r$) and
      the term $r\Psi_{xx}$ can only balance at $w_x = 0$, which kills the
      spatial scaling and with it any weight on $E$. This is the correct and
      physically necessary answer --- the energy eigenvalue is the output of
      the computation, and an algorithm that offered to normalize it away would
      be wrong. It is also the exact analogue of the $p_1 \ldots p_4$ situation
      of \cite{ansatz25}, where zero gauge weight forced a pseudo-division
      argument in place of a symmetry argument.
\item \textbf{Two parameters may be pinned}, taking the count from $10$ to $8$
      and removing $6$ of the $30$ constancy equations.
\item \textbf{Five of the two-parameter pins are defective}, in the precise
      sense of \S\ref{sec:unimodular}:
      \begin{center}
      \begin{tabular}{l l}
      \toprule
      pin & consequence \\
      \midrule
      $a_0 = b_1 = 1$ & $2$-fold cover \\
      $a_0 = c_0 = 1$ & $2$-fold cover \\
      $a_0 = c_1 = 1$ & $3$-fold cover \\
      $a_1 = c_1 = 1$ & $2$-fold cover \\
      $b_0 = c_1 = 1$ & $2$-fold cover \\
      \bottomrule
      \end{tabular}
      \end{center}
      $a_0 = c_0 = 1$ --- normalize the leading and trailing constant
      coefficients, surely the first choice anyone would write down --- doubles
      every component. Of the $55$ pairs, $32$ are clean and $18$ are
      degenerate (rank $<2$), the latter including every pair involving $E$ and
      the equal-weight pairs $\{a_1, b_0\}$ and $\{b_1, c_0\}$.
\end{enumerate}

\section{Tier 2: all linear symmetries}
\label{sec:tier2}

\subsection{The blind spot}

A weight table is the character lattice of a maximal torus, and \emph{a
unipotent element has no character}. Tier~1 is therefore not merely incomplete
but incomplete in a specific, predictable direction, and the consequence has
already been observed empirically. The normalization ladder for ansatz~25 was
designed entirely from a weight table and consequently spent $6$ of the $9$
available group dimensions, missing all three of
\cite{ansatz25}:
\[
  \tau : s \mapsto s + c \;\;(\Ga), \qquad
  \delta : \text{parabolic scaling} \;\;(\Gm), \qquad
  \pi : p \mapsto p + f(t) \;\;(\Ga),
\]
worth three parameters and twelve constancy equations per rung, and --- unlike
the rotation normal form the ladder \emph{did} spend --- each admitting a
complete, exclusion-free chart cover. Two of the three missed generators are
unipotent, which is exactly the shape a character table cannot see.

\subsection{The remedy, and its cost}

The missed directions are nonetheless \emph{linear} on parameter space. The
$\tau$ action, for instance, sends an affine coefficient pair by
\[
  (x_0 + x_1 s) \;\longmapsto\; (x_0 + c\,x_1) + x_1 s ,
  \qquad\text{i.e.}\qquad x_0 \mapsto x_0 + c\,x_1, \;\; x_1 \mapsto x_1 ,
\]
a single Jordan block --- linear, merely not diagonalizable. So widen the
search from diagonal matrices to all of $\gl_N$: seek a vector field
$X \in \gl(P)$ such that

\begin{enumerate}
\item[(i)] $X$ preserves the system: each $X \cdot f$ pseudo-reduces to zero
           against the generators. This is \textbf{linear in the $N^2$ entries
           of $X$}, the reduction supplying the coefficients.
\item[(ii)] $X$ moves the coordinate of interest: $X k \not\equiv 0$ on the
            cell --- an open condition.
\end{enumerate}

One solve then captures torus, unipotent and rotational directions together.
For hydrogen / ansatz~5 that is $11$ parameters, hence $\approx 121$ unknowns
and one reduction per generator per basis element: entirely negligible beside
the $1906$~s of cell~13.

\begin{remark}
Tier~2 is complete for group actions that are \emph{linear on parameter space}.
That covers $\tau$, $\pi$, $\delta$ and the rotations, but it is a modelling
restriction rather than a theorem: parameters may in principle transform
polynomially. Genuine completeness is Tier~3.
\end{remark}

\section{Tier 3: general Lie point symmetries}
\label{sec:tier3}

The complete answer is classical and is already cited in the paper: prolong a
general vector field, impose the infinitesimal criterion that it annihilate the
system modulo the system, and obtain an overdetermined linear PDE system --- the
\emph{determining equations} --- for its coefficients \cite{Olver}. Solving
that system is a differential-elimination problem, so the available software is
Reid's \code{rifsimp}, Maple's \code{PDEtools[Infinitesimals]}, and the Thomas
decomposition of \cite{ThomasDecomp}. The machinery that would compute the
symmetries is the machinery the pipeline already runs.

For the present purpose, however, the more directly relevant strand of Olver's
work is the second one: \emph{moving frames}
\cite{OlverEquivalence,FelsOlverI,FelsOlverII}. The normalization construction
of \S\ref{sec:criterion} below is precisely a Fels--Olver moving frame ---
choose a cross-section to the group orbits, solve the normalization equations
for the group parameters, substitute back --- and the regularity conditions of
that theory are exactly the trichotomy of Criterion~\ref{crit:setone}.

\section{The per-cell criterion: when may one set $k = 1$?}
\label{sec:criterion}

\subsection{Statement}

Consider the moment a decomposition splits, adding $k \neq 0$ to a cell's
inequations. Rather than carry $k$ as a free nonzero parameter, one would like
to pin $k = 1$.

\begin{criterion}[set-it-to-one]
\label{crit:setone}
Let $C$ be a cell whose inequations include $k \neq 0$. Ask whether the
quasi-homogeneity lattice of $C$ contains an integer weight vector $w$ with
$w_k = 1$. There are three cases.
\begin{itemize}
\item[$w_k = 1$:] the $\Gm$ acts by $k \mapsto tk$, so $t = 1/k$ is the
      \emph{unique} group element carrying a given point of $C$ to $\{k=1\}$.
      The slice meets every orbit exactly once. \textbf{Clean pin}: no chart,
      no cover, no residual.
\item[$w_k = d$, $|d| > 1$, and no weight-$1$ vector:] $t^d k = 1$ has $d$
      roots, the slice meets each orbit $d$ times, and the answer is a
      $d$-fold cover. Usable only with an explicit $\mu_d$ quotient
      (\S\ref{sec:unimodular}).
\item[$w_k = 0$ forced:] no scaling moves $k$; it is a genuine modulus and
      cannot be pinned by any Tier-1 argument.
\end{itemize}
\end{criterion}

Two features make this attractive beyond the bare statement.

\emph{The exclusion is free.} A normalization ordinarily buys a parameter at
the price of a chart --- one must handle the excluded locus $k = 0$ separately.
Here the split has \emph{already} excluded it: $k = 0$ is the sibling branch,
which the decomposition was going to compute regardless. The pin spends a
redundancy the algorithm has just handed over at no cost.

\emph{The weight vector reports its own bookkeeping.} Inspect the components of
$w$ on the differential indeterminates and independent variables, as opposed to
those on the parameters and the equation-multiplier. If the former all vanish,
the symmetry is gauge (\S\ref{sec:twokinds}): the solution set is
\emph{literally identical}, and nothing is owed. If not, solutions move, and
the answer at $k = 1$ must be transported over the $\Gm$-orbit --- but the
transport is a substitution applied to the finished answer, never a
recomputation. So the distinction of \S\ref{sec:twokinds} is decided
mechanically, by reading off a sub-vector.

\subsection{Correspondence with moving frames}

Criterion~\ref{crit:setone} is the one-parameter case of the Fels--Olver
regularity trichotomy, and it is worth recording the dictionary, both to place
the construction in the literature and because the general theory then supplies
the multi-parameter version for free:

\begin{center}
\begin{tabular}{l l}
\toprule
Criterion~\ref{crit:setone} & Fels--Olver \cite{FelsOlverI,FelsOlverII} \\
\midrule
$w_k = 1$ achievable      & action \textbf{free} $\Rightarrow$ global cross-section, unique normalization \\
$w_k = d$, $|d| > 1$      & \textbf{locally free}, discrete isotropy $\mu_d$ $\Rightarrow$ local frame, quotient owed \\
$w_k = 0$ forced          & not free in that direction $\Rightarrow$ no cross-section; genuine modulus \\
\bottomrule
\end{tabular}
\end{center}

The pin $k = 1$ is a cross-section; $t = 1/k$ is the moving frame $\rho$; and
the invariants of the reduced problem are the invariantizations
$\iota(F) = F(\rho \cdot z)$.

\subsection{Can the criterion be tested rather than derived?}

One would like to verify the symmetry empirically --- introduce an
indeterminate $\lambda$, apply the scaling, and check that every transformed
generator differential-pseudo-reduces to zero modulo the original system, and
conversely. This is an ideal-equality test by reduction, and it is cheap with
machinery the solver already has.

It is worth doing as a check on a derived $w$, but it is \emph{not} an
alternative route to one. Scaling $k$ alone is essentially never a symmetry; it
closes only when every other symbol is rescaled in the matching way, and
finding that matching companion scaling \emph{is} the weight computation. There
is no test-only shortcut past the derivation --- but the derivation is one
Hermite normal form.

\section{Why per-cell strictly dominates global normalization}
\label{sec:monotone}

The ansatz-25 ladder \cite{ansatz25} normalizes once, globally, before the
decomposition begins. Criterion~\ref{crit:setone} instead fires at each split.
The following shows the latter is never worse and generally better.

\begin{proposition}[monotonicity of available scaling]
\label{prop:monotone}
Let the input system be $w$-quasi-homogeneous. Then every cell of its Thomas
decomposition is $w$-quasi-homogeneous, and the dimension of the scaling group
available on a cell is non-decreasing along any root-to-leaf path.
\end{proposition}

\begin{proof}[Argument]
Three observations.
\begin{enumerate}
\item Every operation the decomposition performs --- pseudo-division,
      formation of initials and separants, resultants, and the differentiations
      that generate prolongations --- preserves weights. Hence
      $w$-quasi-homogeneity is inherited by every cell, and the root lattice is
      available on every cell at no cost.
\item A cell's \emph{new} equations arise as pseudo-remainders of its parents'
      generators, and are therefore themselves $w$-quasi-homogeneous. The rows
      they contribute to $M$ already lie in its row space, so they cannot
      shrink $\ker M$.
\item A cell's equations that \emph{set parameters to zero} delete monomials
      from the surviving generators, hence delete rows from $M$. Deleting rows
      can only enlarge the kernel.
\end{enumerate}
Combining (2) and (3), the kernel is non-decreasing as one descends. \qed
\end{proof}

\begin{remark}
Proposition~\ref{prop:monotone} is the argument for building the criterion into
the split rather than into the preprocessing: deep cells have \emph{at least
as much}, and typically strictly more, normalization available than the root,
and a global up-front pass leaves all of the excess unspent. Cell~13 of
Table~\ref{tab:hydrogen5} is the case in point --- it has two parameters
zeroed, so by (3) its lattice has grown, and it retains seven inequations, so
it offers seven candidate pins. It is also, by an order of magnitude, the cell
whose cost matters.
\end{remark}

\section{What the criterion does not catch}
\label{sec:gap}

Criterion~\ref{crit:setone} is a Tier-1 certificate with a per-cell trigger. It
therefore inherits the blind spot of \S\ref{sec:tier2} in full, and it is worth
being precise about how, because the trigger and the certificate fail
\emph{differently}.

\paragraph{The trigger generalizes; the certificate does not.} The question
``is there a group element carrying every point of this branch to $k=1$?'' is
group-agnostic and is the right question for any $G$. The test
``$\exists\, w \in \ZZ^N$ with $w_k = 1$'' answers it correctly for a torus and
returns \emph{no} for genuine unipotent symmetries that would have worked.

\paragraph{Pure translations are never triggered at all.} A $\Ga$ acting by
$k \mapsto k + c$ has a single orbit on the $k$-line, all of $\CC$. There is no
invariant $\{k \neq 0\}$ locus for a split to produce, so no inequation ever
fires the trigger. The natural normalization is $k = 0$, unconditional, with
neither branch nor chart --- which is precisely the $\pi$ generator of
\cite{ansatz25} (``$p_4 = 0$: lossless, no chart, no split''). This is
\emph{cheaper} than any torus pin, and a criterion keyed to nonzero splits will
never find it. It needs its own trigger.

\paragraph{Unipotents on a plane are triggered on the wrong variable.} For
$\tau : x_0 \mapsto x_0 + c\,x_1$, $x_1 \mapsto x_1$, the orbits are the lines
$\{x_1 = \text{const}\}$ when $x_1 \neq 0$, and fixed points when $x_1 = 0$. So
a branch \emph{does} license a pin --- but the inequation is on $x_1$ while the
coordinate pinned is $x_0$. A criterion that looks for a pin on the variable
the inequation is about is not merely unable to certify this case; it is not
looking in the right place.

\paragraph{The composite case, which is the cautionary one.} In ansatz~25,
$\tau$ composed with the gauge renormalization $\kappa = 1/(1+g_1c)$ acts on
$g_1$ by
\[
   g_1 \;\longmapsto\; \frac{g_1}{1 + c\,g_1} ,
\]
a M\"obius action whose orbits are exactly $\{0\}$ and $\CC^*$. That is
indistinguishable in orbit structure from a scaling, and it is the origin of
the proposed normalization $g_1 \in \{0,1\}$ of \cite{ansatz25}, \S6. So the
trigger \emph{would} have fired correctly on $g_1$ --- and the weight-lattice
certificate would have said \emph{no}, the action being neither diagonal nor
even linear in that coordinate. The trigger was right and the certificate was
wrong. This is the failure mode to design against.

\paragraph{Consequent recommendation.} Tier~1 is attractive because it is a
two-line Hermite normal form. But Tier~2 is also only linear algebra, is
negligible against the cost it is meant to reduce (\S\ref{sec:tier2}), and the
one case in this project where the analysis has been carried out by hand is a
case in which \textbf{Tier~1 alone would have missed everything that was in
fact missed}. The diagnostic should be built at Tier~2 from the start.

\section{What is not a group at all}
\label{sec:notagroup}

A caveat belongs in any write-up of this material, since it bounds what
\emph{any} symmetry algorithm at \emph{any} tier can deliver. Not every
redundancy is a group action.

Two linear ODE operators define the same solution family iff they generate the
same left ideal in the Weyl algebra --- what matters is $D\cdot L$, not $L$.
Restricted to a family with affine coefficients, that relation is a groupoid
with jumping fibres rather than a group action: over a constant-coefficient
operator the fibre $\{(1 + g_1 s)\cdot L\}$ is $1$-dimensional, while over a
generic affine operator it is a point \cite{ansatz25}. The gauge generator
$\kappa$ is only its degree-zero part. No normalization can pin the rest,
because there is no group to quotient by, and a symmetry-detection algorithm
that reported one would be in error.

\section{Positioning with respect to the paper's Conclusion}
\label{sec:positioning}

The Conclusion of \emph{A New Method of Solving PDEs} draws a sharp line
against symmetry methods:

\begin{quote}
\small
\ldots complementary to the \emph{similarity-reduction} methods --- the
classical Lie symmetry method \cite{Olver}, the Clarkson--Kruskal direct method,
and the nonclassical (conditional) symmetries of Bluman and Cole --- which also
reduce a PDE to an ODE without assuming a separable product form, but which are
organized around the (classical or conditional) symmetries of the equation. Our
reductions are instead organized around the differential-algebraic structure of
the ansatz; \textbf{they do not presuppose any symmetry of the PDE}, and they
carry the completeness guarantee of Theorem~1.
\end{quote}

A \S\,Scaling Symmetries subsection reintroduces symmetry, and a referee may
read it as undercutting that claim. The defence is clean and should be stated
explicitly rather than left to inference:

\begin{quote}
The similarity-reduction methods use symmetry \emph{constitutively} --- the
symmetry group is what selects the reduction, so an equation with no symmetry
admits no method. The present use is different in kind: symmetry normalizes the
parameter space of an ansatz that was chosen on entirely other grounds. Delete
the symmetry analysis and the method still runs, still correct and still
complete, merely with more parameters and a slower decomposition. Symmetry here
is an optimization, not a premise, and Theorem~1 is untouched either way.
\end{quote}

One sentence to this effect in the new subsection, with a cross-reference from
the Conclusion, closes the gap.

\section{Recommendations}
\label{sec:recommendations}

\begin{enumerate}
\item \textbf{Adopt the unimodularity check immediately, before any pin.} It is
      a determinant, and \S\ref{sec:hydrogen} shows the natural choice
      $a_0 = c_0 = 1$ is defective in the very ansatz under study. A silent
      $d$-fold cover is the worst class of error here: the run completes, the
      output looks structurally correct, and every component is duplicated.
\item \textbf{Build the diagnostic at Tier~2, not Tier~1} (\S\ref{sec:gap}).
      Cost is comparable; coverage is not.
\item \textbf{Add a separate trigger for translation directions}, which the
      nonzero-split hook cannot reach and which are the cheapest pins available
      (no chart, no split).
\item \textbf{Run the criterion as a read-only diagnostic first.} For each cell,
      report which inequations admit a clean pin, which give a $d$-cover, and
      which are moduli. This measures the payoff --- concentrated, per
      Table~\ref{tab:hydrogen5}, almost entirely in cell~13 --- without
      perturbing the decomposition.
\item \textbf{Retain an unpinned mode.} Pinning inside the decomposition means
      the output is no longer the \emph{canonical} Thomas decomposition of the
      original system. Disjointness and the partition identity survive;
      canonicity and comparability against a reference run do not.
\item \textbf{Consider the dual construction.} Since the Thomas decomposition is
      canonical, a symmetry group \emph{permutes cells}. Where a $\Gm$ carries
      cell $i$ to cell $j$, one does not pin --- one discards $j$ as a duplicate
      and records the transport. Given the cost concentration of
      Table~\ref{tab:hydrogen5}, this may be worth more than pinning. (Cells 1
      and 2, and 3 and 4, return identical prime lists, which is suggestive;
      but the likelier explanation is GTZ absorbing limit strata, and the
      question should be settled before anything is claimed.)
\end{enumerate}

\section{What was checked, and how}
\label{sec:checked}

Following the discipline of \cite{ansatz25}, \S7, the computational claims here
were executed rather than asserted.

\begin{itemize}
\item \textbf{The hydrogen / ansatz 5 weight table} was verified one
      generator at a time by the script
      \code{symmetry-weights-check.py} in this
      directory: all nine ansatz equations, the PDE, and both inequations are
      quasi-homogeneous under the stated weights, each with a single monomial
      weight. The script also enumerates all $55$ two-parameter pins and
      classifies them by $|\det|$, producing the $32$/$5$/$18$ split of
      \S\ref{sec:hydrogen}.
\item \textbf{$w_x = 0$ is forced} by the PDE: the monomials $r\Psi_{xx}$ and
      $\Psi$ carry weights $w_\Psi - w_x$ and $w_\Psi$, which agree only at
      $w_x = 0$; $w_E = -w_x = 0$ follows. Physically, the Coulomb term fixes a
      length scale, so hydrogen at fixed charge has no spatial scaling
      symmetry, and the eigenvalue is correspondingly unnormalizable.
\item \textbf{Cross-validation against \cite{ansatz25}.} The computed
      $\mu$-weights for $(a_0,a_1,b_0,b_1,c_0,c_1)$, namely
      $(2,1,1,0,0,-1)$, reproduce that note's independently derived
      general rule --- $\mu$-weights $w_j$ and $w_j - 1$ for the pair
      multiplying a jet monomial of derivative order $w_j$ --- at
      $w_j = 2, 1, 0$.
\item \textbf{The $2$-fold cover of $\Psi'' + b\Psi' + \Psi$} was checked
      directly: the residual subgroup is $\{t_x = \pm 1\}$ acting by
      $b \mapsto \pm b$, realized by $x \mapsto -x$.
\item \textbf{Timings} in Table~\ref{tab:hydrogen5} are read from
      \code{hydrogen-5-membership.out}, the run of 5--6~September 2026 on
      samsung.
\end{itemize}

\section{Gaps}
\label{sec:gaps}

\begin{itemize}
\item \textbf{No Tier-2 computation has been performed for hydrogen /
      ansatz~5.} The lattice of \S\ref{sec:hydrogen} is the torus only. Whether
      that system admits unipotent directions analogous to ansatz~25's $\tau$
      is open; the inner variable $v = v_1x + v_2y + v_3z + v_4r$ carries no
      constant term, so the obvious $\tau$-analogue is \emph{not} available,
      but this is an observation, not a proof of absence.
\item \textbf{The payoff is projected, not measured.} No run has been made with
      any pin. Two parameters and six constancy equations is exact arithmetic;
      that this translates into a shorter cell~13 is a hypothesis. The same
      caution applies as in \cite{ansatz25}, \S6, where it was noted that the
      cost driver might be the structure of a stratum rather than raw parameter
      count --- and cell~13, where the ODE degenerates in order, is exactly
      such a structural stratum.
\item \textbf{Pinning a polynomial rather than a parameter} is unimplemented.
      Cell~13's seven inequations are in general polynomials, not bare
      parameters. By Proposition~\ref{prop:monotone} they are $w$-homogeneous,
      so the criterion applies verbatim; but pinning ``this polynomial $= 1$''
      is a coordinate change rather than a substitution, and is more work.
\item \textbf{Discrete symmetries are out of scope} at every tier. Tiers 1--3
      compute connected groups; the reflection $x \mapsto -x$ of
      \S\ref{sec:unimodular} appeared only as a residual isotropy, and nothing
      here would have found it \emph{ab initio}.
\item \textbf{The consistency-locus half is unfinished.} The companion
      \code{--locus consistency} run for hydrogen / ansatz~5 was killed at
      $1$h$06$m and $23.7$~GB, still inside the Thomas decomposition. The cost
      figures of \S\ref{sec:cost} are for the membership locus only.
\end{itemize}

\begin{thebibliography}{99}

\bibitem{ansatz25}
\emph{The symmetry group of ansatz 25}, companion note, 2~September 2026,
\nolinkurl{~/Papers/NewMethod/ansatz-25-symmetry-group.md}. [Lie-group analysis
of the ansatz-25 parameter space: the $\Gamma \times H$ splitting, the weight
table, the unimodularity criterion, the three unspent generators $\tau$,
$\delta$, $\pi$, and the Weyl-algebra groupoid of \S\ref{sec:notagroup}.]

\bibitem{Olver}
P.~J. Olver, \emph{Applications of Lie Groups to Differential Equations}, second
ed., Graduate Texts in Mathematics, vol.~107, Springer, New York, 1993.
\url{https://doi.org/10.1007/978-1-4612-4350-2}. [Ch.~2: prolongation formula,
infinitesimal criterion, determining equations --- Tier~3.]

\bibitem{OlverEquivalence}
P.~J. Olver, \emph{Equivalence, Invariants and Symmetry}, Cambridge University
Press, 1995.

\bibitem{FelsOlverI}
M. Fels, P.~J. Olver, Moving coframes~I: a practical algorithm, \emph{Acta Appl.
Math.} \textbf{51} (1998) 161--213.

\bibitem{FelsOlverII}
M. Fels, P.~J. Olver, Moving coframes~II: regularization and theoretical
foundations, \emph{Acta Appl. Math.} \textbf{55} (1999) 127--208.
[Freeness / local freeness / regularity --- the trichotomy of
Criterion~\ref{crit:setone}.]

\bibitem{HubertLabahnISSAC}
\'E. Hubert, G. Labahn, Rational invariants of scalings from Hermite normal
forms, in: \emph{Proc.\ ISSAC 2012}, ACM, pp.~219--226.

\bibitem{HubertLabahnFoCM}
\'E. Hubert, G. Labahn, Scaling invariants and symmetry reduction of dynamical
systems, \emph{Found. Comput. Math.} \textbf{13} (2013) 479--516.
[The Hermite-normal-form algorithm of \S\ref{sec:tier1}.]

\bibitem{ThomasDecomp}
T. B\"achler, V.~P. Gerdt, M. Lange-Hegermann, D. Robertz, Algorithmic Thomas
decomposition of algebraic and differential systems, \emph{J. Symbolic Comput.}
\textbf{47}(10) (2012) 1233--1266.
\url{https://doi.org/10.1016/j.jsc.2011.12.043}.

\bibitem{hubert}
\'E. Hubert, Notes on triangular sets and triangulation-decomposition
algorithms~I: Polynomial systems, LNCS vol.~2630, Springer, 2003, pp.~1--39.

\bibitem{hubert2}
\'E. Hubert, Notes on triangular sets and triangulation-decomposition
algorithms~II: Differential systems, LNCS vol.~2630, Springer, 2003,
pp.~40--87.

\end{thebibliography}

\bigskip
\noindent\emph{Prepared as computational and literature backing for the
\S\,Scaling Symmetries subsection of} A New Method of Solving PDEs. \emph{The
Tier-1 lattice of \S\ref{sec:hydrogen} and the pin classification were computed
and verified live; see \S\ref{sec:checked}. Written by Claude (Opus~5) in an
interactive session of 6~September 2026 --- reviewed and edited by the author;
see} \code{AI-CONTRIBUTIONS.md}.

\end{document}
